% ============================================================
% Question Bank: ch05-03 Solving Equations with the Distributive Property
% Topic Title: Solving Equations with the Distributive Property
% CCSS: 7.EE.B.4.A
% Questions: 30 (18 MC + 7 SA + 5 Graphical)
% ============================================================

% --- Multiple Choice Questions (18) ---

\begin{practiceQuestion}{5-3-q01}{mc}
\begin{questionText}
Solve $2(x + 6) = 20$.
\end{questionText}
\choiceA{$x = 4$}
\choiceB{$x = 7$}
\choiceC{$x = 10$}
\choiceD{$x = 14$}
\correctAnswer{A}
\explanation{Divide both sides by $2$: $x + 6 = 10$. Subtract $6$: $x = 4$. Check: $2(4 + 6) = 2(10) = 20$ \checkmark.}
\end{practiceQuestion}

\begin{practiceQuestion}{5-3-q02}{mc}
\begin{questionText}
Solve $3(n - 4) = 15$.
\end{questionText}
\choiceA{$n = 9$}
\choiceB{$n = 3$}
\choiceC{$n = \dfrac{19}{3}$}
\choiceD{$n = 1$}
\correctAnswer{A}
\explanation{Divide by $3$: $n - 4 = 5$. Add $4$: $n = 9$. Check: $3(9 - 4) = 3(5) = 15$ \checkmark.}
\end{practiceQuestion}

\begin{practiceQuestion}{5-3-q03}{mc}
\begin{questionText}
What is the first step to solve $5(x + 2) = 35$ using the ``divide first'' method?
\end{questionText}
\choiceA{Subtract $2$ from both sides}
\choiceB{Distribute $5$ to get $5x + 10 = 35$}
\choiceC{Divide both sides by $5$}
\choiceD{Subtract $5$ from both sides}
\correctAnswer{C}
\explanation{The ``divide first'' method divides both sides by $5$: $x + 2 = 7$.}
\end{practiceQuestion}

\begin{practiceQuestion}{5-3-q04}{mc}
\begin{questionText}
Solve $-4(y + 3) = 28$.
\end{questionText}
\choiceA{$y = 4$}
\choiceB{$y = -10$}
\choiceC{$y = 10$}
\choiceD{$y = -4$}
\correctAnswer{B}
\explanation{Divide by $-4$: $y + 3 = -7$. Subtract $3$: $y = -10$. Check: $-4(-10 + 3) = -4(-7) = 28$ \checkmark.}
\end{practiceQuestion}

\begin{practiceQuestion}{5-3-q05}{mc}
\begin{questionText}
Solve $6(x - 1) = -18$.
\end{questionText}
\choiceA{$x = -2$}
\choiceB{$x = 2$}
\choiceC{$x = -4$}
\choiceD{$x = 4$}
\correctAnswer{A}
\explanation{Divide by $6$: $x - 1 = -3$. Add $1$: $x = -2$. Check: $6(-2 - 1) = 6(-3) = -18$ \checkmark.}
\end{practiceQuestion}

\begin{practiceQuestion}{5-3-q06}{mc}
\begin{questionText}
Which equation is equivalent to $4(x + 5) = 36$?
\end{questionText}
\choiceA{$4x + 5 = 36$}
\choiceB{$4x + 20 = 36$}
\choiceC{$4x + 9 = 36$}
\choiceD{$x + 20 = 36$}
\correctAnswer{B}
\explanation{Distribute $4$: $4 \times x + 4 \times 5 = 4x + 20$. So the equation becomes $4x + 20 = 36$.}
\end{practiceQuestion}

\begin{practiceQuestion}{5-3-q07}{mc}
\begin{questionText}
Solve $3(2x + 1) = 21$.
\end{questionText}
\choiceA{$x = 3$}
\choiceB{$x = 6$}
\choiceC{$x = 10$}
\choiceD{$x = \dfrac{20}{6}$}
\correctAnswer{A}
\explanation{Distribute: $6x + 3 = 21$. Subtract $3$: $6x = 18$. Divide by $6$: $x = 3$. Check: $3(2 \cdot 3 + 1) = 3(7) = 21$ \checkmark.}
\end{practiceQuestion}

\begin{practiceQuestion}{5-3-q08}{mc}
\begin{questionText}
Solve $-2(3x - 5) = 16$.
\end{questionText}
\choiceA{$x = -1$}
\choiceB{$x = 1$}
\choiceC{$x = -\dfrac{13}{3}$}
\choiceD{$x = 3$}
\correctAnswer{A}
\explanation{Distribute: $-6x + 10 = 16$. Subtract $10$: $-6x = 6$. Divide by $-6$: $x = -1$. Check: $-2(3(-1) - 5) = -2(-8) = 16$ \checkmark.}
\end{practiceQuestion}

\begin{practiceQuestion}{5-3-q09}{mc}
\begin{questionText}
Each side of a square is $(x + 7)$ cm. The perimeter is $60$ cm. What is $x$?
\end{questionText}
\choiceA{$x = 15$}
\choiceB{$x = 53$}
\choiceC{$x = 8$}
\choiceD{$x = 22$}
\correctAnswer{C}
\explanation{Perimeter $= 4(x + 7) = 60$. Divide by $4$: $x + 7 = 15$. Subtract $7$: $x = 8$.}
\end{practiceQuestion}

\begin{practiceQuestion}{5-3-q10}{mc}
\begin{questionText}
Solve $5(n + 3) = -10$.
\end{questionText}
\choiceA{$n = -5$}
\choiceB{$n = 1$}
\choiceC{$n = -1$}
\choiceD{$n = -\dfrac{25}{5}$}
\correctAnswer{A}
\explanation{Divide by $5$: $n + 3 = -2$. Subtract $3$: $n = -5$. Check: $5(-5 + 3) = 5(-2) = -10$ \checkmark.}
\end{practiceQuestion}

\begin{practiceQuestion}{5-3-q11}{mc}
\begin{questionText}
A student solved $3(x + 4) = 30$ and wrote: $3x + 4 = 30$, $3x = 26$, $x = \frac{26}{3}$. What was the student's error?
\end{questionText}
\choiceA{Divided by $3$ instead of subtracting}
\choiceB{Did not distribute $3$ to the $4$}
\choiceC{Subtracted $4$ from the wrong side}
\choiceD{Forgot to check the answer}
\correctAnswer{B}
\explanation{The student wrote $3x + 4$ instead of $3x + 12$. Distributing correctly gives $3x + 12 = 30$.}
\end{practiceQuestion}

\begin{practiceQuestion}{5-3-q12}{mc}
\begin{questionText}
Solve $7(x - 2) = 0$.
\end{questionText}
\choiceA{$x = 0$}
\choiceB{$x = -2$}
\choiceC{$x = 2$}
\choiceD{$x = 7$}
\correctAnswer{C}
\explanation{Divide by $7$: $x - 2 = 0$. Add $2$: $x = 2$. Check: $7(2 - 2) = 7(0) = 0$ \checkmark.}
\end{practiceQuestion}

\begin{practiceQuestion}{5-3-q13}{mc}
\begin{questionText}
Solve $-5(x + 1) = -30$.
\end{questionText}
\choiceA{$x = 5$}
\choiceB{$x = -7$}
\choiceC{$x = 7$}
\choiceD{$x = -5$}
\correctAnswer{A}
\explanation{Divide by $-5$: $x + 1 = 6$. Subtract $1$: $x = 5$. Check: $-5(5 + 1) = -5(6) = -30$ \checkmark.}
\end{practiceQuestion}

\begin{practiceQuestion}{5-3-q14}{mc}
\begin{questionText}
Solve $2(4x - 3) = 26$.
\end{questionText}
\choiceA{$x = 3$}
\choiceB{$x = 4$}
\choiceC{$x = 2$}
\choiceD{$x = 5$}
\correctAnswer{B}
\explanation{Distribute: $8x - 6 = 26$. Add $6$: $8x = 32$. Divide by $8$: $x = 4$. Check: $2(4 \cdot 4 - 3) = 2(13) = 26$ \checkmark.}
\end{practiceQuestion}

\begin{practiceQuestion}{5-3-q15}{mc}
\begin{questionText}
Three friends split the cost of a gift equally. Each person also pays $\$4$ for wrapping. Each person pays $\$16$ total. What is the cost $c$ of the gift?
\end{questionText}
\choiceA{$c = 48$}
\choiceB{$c = 36$}
\choiceC{$c = 12$}
\choiceD{$c = 60$}
\correctAnswer{B}
\explanation{Each pays $\frac{c}{3} + 4 = 16$. This means $\frac{c}{3} = 12$, so $c = 36$. Or: total paid $= 3 \times 16 = 48$, wrapping $= 3 \times 4 = 12$, gift $= 48 - 12 = 36$.}
\end{practiceQuestion}

\begin{practiceQuestion}{5-3-q16}{mc}
\begin{questionText}
Solve $8(x + 1) = 48$.
\end{questionText}
\choiceA{$x = 5$}
\choiceB{$x = 6$}
\choiceC{$x = 7$}
\choiceD{$x = \dfrac{47}{8}$}
\correctAnswer{A}
\explanation{Divide by $8$: $x + 1 = 6$. Subtract $1$: $x = 5$. Check: $8(5 + 1) = 8(6) = 48$ \checkmark.}
\end{practiceQuestion}

\begin{practiceQuestion}{5-3-q17}{mc}
\begin{questionText}
Solve $-3(2n + 4) = -24$.
\end{questionText}
\choiceA{$n = 6$}
\choiceB{$n = -2$}
\choiceC{$n = 2$}
\choiceD{$n = 4$}
\correctAnswer{C}
\explanation{Divide by $-3$: $2n + 4 = 8$. Subtract $4$: $2n = 4$. Divide by $2$: $n = 2$. Check: $-3(2(2)+4) = -3(8) = -24$ \checkmark.}
\end{practiceQuestion}

\begin{practiceQuestion}{5-3-q18}{mc}
\begin{questionText}
Solve $10(x - 5) = 30$.
\end{questionText}
\choiceA{$x = 8$}
\choiceB{$x = 3$}
\choiceC{$x = -2$}
\choiceD{$x = 35$}
\correctAnswer{A}
\explanation{Divide by $10$: $x - 5 = 3$. Add $5$: $x = 8$. Check: $10(8 - 5) = 10(3) = 30$ \checkmark.}
\end{practiceQuestion}

% --- Short Answer Questions (7) ---

\begin{practiceQuestion}{5-3-q19}{sa}
\begin{questionText}
Solve $4(x + 9) = 52$.
\end{questionText}
\correctAnswer{$x = 4$}
\explanation{Divide by $4$: $x + 9 = 13$. Subtract $9$: $x = 4$. Check: $4(4 + 9) = 4(13) = 52$ \checkmark.}
\end{practiceQuestion}

\begin{practiceQuestion}{5-3-q20}{sa}
\begin{questionText}
Solve $-6(n - 2) = 42$.
\end{questionText}
\correctAnswer{$n = -5$}
\explanation{Divide by $-6$: $n - 2 = -7$. Add $2$: $n = -5$. Check: $-6(-5 - 2) = -6(-7) = 42$ \checkmark.}
\end{practiceQuestion}

\begin{practiceQuestion}{5-3-q21}{sa}
\begin{questionText}
Solve $5(2x - 3) = 35$.
\end{questionText}
\correctAnswer{$x = 5$}
\explanation{Divide by $5$: $2x - 3 = 7$. Add $3$: $2x = 10$. Divide by $2$: $x = 5$. Check: $5(2(5) - 3) = 5(7) = 35$ \checkmark.}
\end{practiceQuestion}

\begin{practiceQuestion}{5-3-q22}{sa}
\begin{questionText}
Solve $2(x + 7) = -6$.
\end{questionText}
\correctAnswer{$x = -10$}
\explanation{Divide by $2$: $x + 7 = -3$. Subtract $7$: $x = -10$. Check: $2(-10 + 7) = 2(-3) = -6$ \checkmark.}
\end{practiceQuestion}

\begin{practiceQuestion}{5-3-q23}{sa}
\begin{questionText}
Each side of an equilateral triangle is $(2x + 5)$ cm. The perimeter is $57$ cm. Find $x$.
\end{questionText}
\correctAnswer{$x = 7$}
\explanation{$3(2x + 5) = 57$. Divide by $3$: $2x + 5 = 19$. Subtract $5$: $2x = 14$. Divide by $2$: $x = 7$.}
\end{practiceQuestion}

\begin{practiceQuestion}{5-3-q24}{sa}
\begin{questionText}
Solve $-8(x + 2) = 24$.
\end{questionText}
\correctAnswer{$x = -5$}
\explanation{Divide by $-8$: $x + 2 = -3$. Subtract $2$: $x = -5$. Check: $-8(-5 + 2) = -8(-3) = 24$ \checkmark.}
\end{practiceQuestion}

\begin{practiceQuestion}{5-3-q25}{sa}
\begin{questionText}
Solve $9(x - 3) = 45$.
\end{questionText}
\correctAnswer{$x = 8$}
\explanation{Divide by $9$: $x - 3 = 5$. Add $3$: $x = 8$. Check: $9(8 - 3) = 9(5) = 45$ \checkmark.}
\end{practiceQuestion}

% --- Graphical Questions (3 GMC + 2 GSA) ---

\begin{practiceQuestion}{5-3-q26}{gmc}
\begin{questionText}
The rectangle below has an area of $56$ square cm. What is the value of $x$?

\begin{center}
\begin{tikzpicture}[scale=0.7]
  \draw[thick,fill=funBlue!10] (0,0) rectangle (7,3);
  \draw[<->] (0,-0.5) -- (7,-0.5);
  \node[below] at (3.5,-0.5) {$7$ cm};
  \draw[<->] (-0.5,0) -- (-0.5,3);
  \node[left] at (-0.5,1.5) {$(x + 3)$ cm};
\end{tikzpicture}
\end{center}
\end{questionText}
\choiceA{$x = 5$}
\choiceB{$x = 8$}
\choiceC{$x = 11$}
\choiceD{$x = 4$}
\correctAnswer{A}
\explanation{Area $= 7(x + 3) = 56$. Divide by $7$: $x + 3 = 8$. Subtract $3$: $x = 5$.}
\end{practiceQuestion}

\begin{practiceQuestion}{5-3-q27}{gmc}
\begin{questionText}
The diagram shows two equal groups. Which equation represents this, and what is $x$?

\begin{center}
\begin{tikzpicture}[scale=0.6]
  % Group 1
  \draw[thick,fill=funBlue!20,rounded corners] (-0.5,-0.5) rectangle (3.5,2.5);
  \node[draw,fill=funBlue!40,minimum size=0.6cm] at (0.5,1.5) {$x$};
  \node at (1.5,1.5) {$+$};
  \node[draw,fill=funOrange!40,minimum size=0.6cm] at (2.5,1.5) {$4$};
  \node at (1.5,0.2) {Group 1};
  % Group 2
  \draw[thick,fill=funBlue!20,rounded corners] (5,-0.5) rectangle (9,2.5);
  \node[draw,fill=funBlue!40,minimum size=0.6cm] at (6,1.5) {$x$};
  \node at (7,1.5) {$+$};
  \node[draw,fill=funOrange!40,minimum size=0.6cm] at (8,1.5) {$4$};
  \node at (7,0.2) {Group 2};
  % Total
  \draw[decorate,decoration={brace,amplitude=8pt}] (9.2,2.7) -- (9.2,-0.7);
  \node[right] at (9.5,1) {$= 22$};
\end{tikzpicture}
\end{center}
\end{questionText}
\choiceA{$2(x + 4) = 22$, $x = 7$}
\choiceB{$2x + 4 = 22$, $x = 9$}
\choiceC{$x + 8 = 22$, $x = 14$}
\choiceD{$2(x + 4) = 22$, $x = 3$}
\correctAnswer{A}
\explanation{Two identical groups of $(x + 4)$: $2(x + 4) = 22$. Divide by $2$: $x + 4 = 11$. Subtract $4$: $x = 7$.}
\end{practiceQuestion}

\begin{practiceQuestion}{5-3-q28}{gmc}
\begin{questionText}
Look at the number line. The expression $3(x + 2)$ equals the distance from $A$ to $B$. What is $x$?

\begin{center}
\begin{tikzpicture}[scale=0.6]
  \draw[thick,->] (-1,0) -- (16,0);
  \foreach \x in {0,3,...,15} \draw (\x,0.15) -- (\x,-0.15) node[below] {$\x$};
  \fill[funBlue] (0,0) circle (4pt) node[above=4pt] {$A$};
  \fill[funOrange] (15,0) circle (4pt) node[above=4pt] {$B$};
  \draw[<->,thick,funBlue] (0,0.6) -- (15,0.6);
  \node[above] at (7.5,0.6) {$3(x+2)$};
\end{tikzpicture}
\end{center}
\end{questionText}
\choiceA{$x = 5$}
\choiceB{$x = 7$}
\choiceC{$x = 3$}
\choiceD{$x = 13$}
\correctAnswer{C}
\explanation{Distance from $A$ to $B$ is $15$. $3(x + 2) = 15$. Divide by $3$: $x + 2 = 5$. Subtract $2$: $x = 3$.}
\end{practiceQuestion}

\begin{practiceQuestion}{5-3-q29}{gsa}
\begin{questionText}
The total area of the figure below is $84$ square cm. Find $x$.

\begin{center}
\begin{tikzpicture}[scale=0.5]
  % 6 identical rectangles side by side
  \foreach \i in {0,...,5} {
    \draw[thick,fill=funBlue!15] (\i*2,0) rectangle (\i*2+2,3);
  }
  \draw[<->] (0,-0.7) -- (12,-0.7);
  \node[below] at (6,-0.7) {$6$ rectangles, each width $(x + 1)$ cm};
  \draw[<->] (-0.7,0) -- (-0.7,3);
  \node[left] at (-0.7,1.5) {$2$ cm};
  \node at (6,4.2) {Total area $= 84$ sq. cm};
\end{tikzpicture}
\end{center}
\end{questionText}
\correctAnswer{$x = 6$}
\explanation{Total area $= 6 \times 2 \times (x + 1) = 84$. So $12(x + 1) = 84$. Divide by $12$: $x + 1 = 7$. Subtract $1$: $x = 6$.}
\end{practiceQuestion}

\begin{practiceQuestion}{5-3-q30}{gsa}
\begin{questionText}
Four friends go to a museum. Each person pays for a ticket and a $\$3$ audio guide. The total for all four people is $\$68$. The equation $4(t + 3) = 68$ represents this situation, where $t$ is the ticket price. Use the model below to find $t$.

\begin{center}
\begin{tikzpicture}[scale=0.6]
  % 4 bars
  \foreach \i in {0,...,3} {
    \draw[thick,fill=funBlue!20] (0,\i*1.2) rectangle (6,\i*1.2+0.9);
    \draw[thick,fill=funOrange!30] (6,\i*1.2) rectangle (8,\i*1.2+0.9);
    \node at (3,\i*1.2+0.45) {$t$};
    \node at (7,\i*1.2+0.45) {$\$3$};
  }
  \node[left] at (-0.3,1.8) {$4$ people};
  \draw[decorate,decoration={brace,amplitude=6pt}] (8.3,4.5) -- (8.3,0);
  \node[right] at (8.7,2.25) {$\$68$};
\end{tikzpicture}
\end{center}
\end{questionText}
\correctAnswer{$t = \$14$}
\explanation{$4(t + 3) = 68$. Divide by $4$: $t + 3 = 17$. Subtract $3$: $t = 14$. Each ticket costs $\$14$.}
\end{practiceQuestion}
